% \section{Vấn đề nghiên cứu}
% Hệ thống văn bản quy phạm pháp luật Việt Nam được cấu trúc phân tầng chặt chẽ từ Luật, Nghị định đến Thông tư, kéo theo mạng lưới quan hệ sửa đổi, bổ sung và bãi bỏ phức tạp với hơn một nửa số văn bản bị điều chỉnh theo thời gian. Các nghiên cứu RAG truyền thống cùng mô hình đồ thị tri thức tĩnh tiêu biểu như SBV-LawGraph \cite{sbvlawgraph} hiện chủ yếu tiếp cận bài toán qua pipeline tuyến tính khép kín từ tiếp nhận truy vấn, gom cụm lân cận một bước đến sinh câu trả lời. 

% Mô hình vận hành tĩnh này bộc lộ ba nút thắt kỹ thuật lớn. Thứ nhất là hiện tượng bùng nổ độ dài prompt và suy giảm độ chính xác ngữ cảnh. Do một văn bản sửa đổi thường phân tán tác động lên nhiều điều khoản độc lập, việc thu thập cơ học toàn bộ các nút liên kết khiến hơn 60\% dữ liệu đưa vào mô hình là thông tin nhiễu, đẩy giá trị Precision@2 xuống ngưỡng 0.37--0.39 \cite{sbvlawgraph} và trực tiếp gây ra lỗi bịa đặt căn cứ pháp lý \cite{lexagenthallu}. Thứ hai, hệ thống tĩnh hoàn toàn thiếu năng lực nhận thức dòng thời gian \cite{fan2026temporal}, dẫn đến việc không thể xác định điều khoản nào thực sự còn hiệu lực khi xảy ra các chuỗi sửa đổi bắc cầu qua nhiều thời kỳ. Cuối cùng, pipeline cố định không thể tự lượng giá mức độ đầy đủ của ngữ cảnh, thiếu khả năng phân rã bài toán đa chặng, không thể quay lui khi gặp nhánh cụt và không có cơ chế chủ động đối thoại để làm rõ dữ kiện còn thiếu.

% \section{Đối tượng thụ hưởng và kịch bản ứng dụng}
% Hệ thống hướng tới việc hỗ trợ chuyên viên pháp chế doanh nghiệp, luật sư và đội ngũ kiểm định tuân thủ trong quá trình rà soát tính hợp pháp của các giao dịch thực tế. Năng lực của hệ thống được kiểm chứng qua ba kịch bản vận hành chủ đạo: truy vết chuỗi văn bản sửa đổi, thẩm định hồ sơ kết hợp làm rõ thông tin và xác thực hiệu lực trực tuyến gắn với phòng thủ chỉ thị ngầm.

% Trong kịch bản truy vết, hệ thống tự động bám theo đồ thị quan hệ từ Luật gốc qua các Nghị định và Thông tư điều chỉnh để trích xuất quy chuẩn áp dụng hiện hành, đồng thời loại bỏ triệt để các điều khoản đã bị bãi bỏ. Ở khâu thẩm định hồ sơ, khi tiếp nhận hợp đồng thiếu các tham số pháp lý tiên quyết như hình thức tổ chức doanh nghiệp để xác định trần sở hữu, tác tử không tự đưa ra giả định ngầm mà tạm dừng tiến trình để yêu cầu người dùng bổ sung dữ liệu. Cuối cùng, hệ thống đối soát trực tiếp số hiệu văn bản cũ với Cơ sở dữ liệu Quốc gia về Văn bản pháp luật \cite{fan2026temporal} nhằm cảnh báo quy định thay thế, kết hợp với các bộ lọc tiền xử lý để phát hiện và vô hiệu hóa các chỉ thị ngầm chèn trong tài liệu \cite{greshake2023promptinjection}.

% \section{Luận điểm lựa chọn kiến trúc Agent thay cho Workflow}
% Một luồng công việc với các nhánh điều kiện viết sẵn hoàn toàn bất khả thi trong miền tri thức luật học bởi không gian tìm kiếm không thể tiền định hóa. Đường dẫn tra cứu pháp lý biến động theo từng vụ việc: một số truy vấn kết thúc ngay tại Thông tư chuyên ngành, trong khi những tình huống phức tạp lại đòi hỏi mở rộng liên ngành sang Luật Doanh nghiệp hoặc Bộ luật Dân sự. Việc cố gắng bao quát hàng nghìn kịch bản tra cứu bằng cấu trúc rẽ nhánh tĩnh sẽ dẫn đến tình trạng bùng nổ tổ hợp mã nguồn. Ngược lại, cơ chế tác tử với năng lực lập kế hoạch động và điều phối công cụ linh hoạt cho phép hệ thống tự lượng giá trạng thái thông tin tại từng bước để quyết định tiếp tục truy vết hay dừng lại.

% Bên cạnh đó, quy trình tuyến tính một chiều của workflow không có cơ chế tự phát hiện và cô lập ảo giác. Bằng việc tích hợp các kỹ thuật Self-RAG \cite{asai2024selfrag}, Reflexion \cite{shinn2023reflexion}, FActScore \cite{min2023factscore} và GANDR \cite{qian2026gandr}, hệ thống thiết lập một chu trình phản biện khép kín: nội dung sinh ra từ tác tử soạn thảo được tác tử kiểm toán bóc tách thành các mệnh đề nguyên tử để đối soát với văn bản gốc. Bất kỳ sai lệch trích dẫn hoặc viện dẫn điều khoản hết hiệu lực nào cũng sẽ kích hoạt lượt truy xuất mới nhằm hiệu chỉnh kết quả. Đồng thời, trước các trạng thái mơ hồ vốn rất phổ biến do yếu tố thời gian và tư cách chủ thể quy định \cite{fan2026temporal}, tác tử có thể chủ động ngắt tiến trình để tương tác với con người, tránh các phán đoán sai lầm mang tính hệ thống của các quy trình cố định.

% \section{Các trụ cột kỹ thuật trọng tâm}
% Nghiên cứu tập trung triển khai bốn trục giải pháp kỹ thuật có tính liên kết chặt chẽ. Trọng tâm của hệ thống là kiến trúc đa tác tử dưới sự chỉ huy của bộ điều phối trung tâm, phân chia vai trò độc lập giữa tác tử định tuyến đồ thị \cite{sbvlawgraph}, tác tử tra cứu thời gian thực trên cổng thông tin điện tử \texttt{vbpl.vn} \cite{fan2026temporal}, tác tử soạn thảo kết luận tuân thủ và tác tử kiểm định mệnh đề \cite{qian2026gandr, min2023factscore}. 

% Vận hành trên cấu trúc này là chính sách truy xuất có định hướng, thay thế cơ chế gom cụm một bước bằng chiến lược tìm kiếm đa chặng chỉ men theo các quan hệ pháp lý mang tính ràng buộc như sửa đổi, thay thế hoặc hướng dẫn thi hành. Tính tin cậy của văn bản sinh ra được bảo đảm thông qua vòng lặp tự phản biện và kiểm chứng chéo \cite{asai2024selfrag, shinn2023reflexion}, ràng buộc mọi lập luận đều phải truy nguyên về câu chữ trong văn bản luật hiện hành. Sau cùng, yếu tố an toàn hệ thống được kiểm soát thông qua cơ chế phân quyền công cụ theo giao thức MCP ba cấp độ, tích hợp điểm phê duyệt của con người trước các tác vụ rủi ro và áp dụng kỹ thuật làm sạch dữ liệu đầu vào nhằm ngăn ngừa nguy cơ tấn công chỉ thị gián tiếp \cite{greshake2023promptinjection}.